\pdfmapfile{+symbols.map}
\pdfmapfile{+cm.map}
\pdfmapfile{+cmextra.map}
\documentclass[11pt]{article}
\usepackage[a4paper,margin=28mm]{geometry}
\usepackage{amsmath,amssymb,amsthm}
\usepackage{microtype}
\usepackage[colorlinks=true,urlcolor=blue]{hyperref}
\setlength{\emergencystretch}{2em}
\newtheorem{proposition}{Proposition}
\title{A translation counterexample to TLMC 00000002685}
\author{Galaxy-0}
\date{12 October 2026}
\begin{document}
\maketitle
\begin{abstract}
The stated orbit-closure dichotomy fails for a single polynomial automorphism
of the affine plane. Translation in one coordinate has constant Jacobian
determinant one, no preperiodic points, and an orbit contained in a proper
closed line. We give a characteristic-zero argument and a kernel-checked
Lean formalization on the rational affine plane. The formalization uses
actual polynomial zero sets, not an axiomatized closure operation.
\end{abstract}

\section{Statement and interpretation}
TLMC 00000002685 proposes that an orbit closure for a multidimensional
endomorphism semigroup is either the preperiodic point set or Zariski dense,
with emptiness of a multiplicative zero set of Jacobian determinants as a
criterion, computably verifiable for at most two maps~\cite{tlmc}.
The text does not restrict the base field, impose degree at least two,
or exclude invariant proper subvarieties.

We use the standard interpretation of endomorphisms as algebraic self-maps
of an affine variety. For a single generator $F$, a point $p$ is preperiodic
if $F^m(p)=F^n(p)$ for some integers $0\leq m<n$. This agrees with finiteness
of its forward orbit. Our example also refutes the weaker reading that the
closure merely consists of preperiodic points. It is not a claim about
endomorphisms required to preserve an algebraic-group identity; translation
would not be a group homomorphism under that different interpretation.

\section{Counterexample}
\begin{proposition}
Let $k$ be any field of characteristic zero. On $\mathbb A^2_k$, the
one-generator semigroup generated by
\[
 F(x,y)=(x+1,y)
\]
has an orbit closure which is neither the preperiodic point set nor Zariski
dense, although its Jacobian determinant and all finite products of that
determinant are everywhere nonzero.
\end{proposition}
\begin{proof}
The map $F$ is a polynomial automorphism, with inverse $(x,y)\mapsto(x-1,y)$.
Induction gives
\[
 F^n(a,b)=(a+n\cdot 1_k,b)\qquad(n\geq0).
\]
If $F^m(a,b)=F^n(a,b)$, then $(n-m)\cdot1_k=0$. Characteristic zero implies
$m=n$. Thus there are no preperiodic points, over $k$ or over any extension
field of $k$.

Take $p=(0,0)$. Its forward orbit is
\[
 \mathcal O(p)=\{(n\cdot1_k,0):n\geq0\}.
\]
Its Zariski closure $C$ contains $p$, and hence is nonempty. The polynomial
$Y$ vanishes on the whole orbit, so $C\subseteq V(Y)$. This is a proper
closed subset, since $(0,1)\notin V(Y)$. Therefore $C$ is not dense, and its
nonemptiness prevents it from equalling the empty preperiodic point set.
Indeed, it does not even consist of preperiodic points.

Finally,
\[
 DF=\begin{pmatrix}1&0\\0&1\end{pmatrix},\qquad \det DF=1.
\]
Every finite product of this determinant is $1$, so its zero set is empty.
The same holds for determinants of iterates, and for compositions or
pullbacks of these constant determinant functions. Thus the proposed
Jacobian condition does not rescue the dichotomy. One generator is already
within the stated limit of at most two maps.
\end{proof}

Only the containment $C\subseteq V(Y)$ is needed. No identification of $C$
with the entire line is used. If a semigroup convention excludes the
identity iterate, the orbit starts at $(1,0)$ instead; it is still nonempty,
contained in the same proper closed line, and contains no preperiodic point.
The same counterexample works on $\mathbb A^d_k$ for every $d\geq2$ by
translating the first coordinate and fixing the others.

\section{What is formally verified}
The accompanying \texttt{Main.lean} is self-contained apart from Lean's
standard distribution. It imports \texttt{Std} and rational arithmetic
lemmas, and uses Lean 4.31.0. Its domain is exactly
$\mathbb Q\times\mathbb Q$. It does not formalize arbitrary fields,
complex points, schemes, or the extension-of-scalars argument above.
The preceding characteristic-zero proof is an ordinary mathematical proof;
only its rational-point instance is kernel-checked here.

A polynomial is represented by a finite expression generated by rational
constants, variables $X,Y$, addition, and multiplication. These expressions
represent all elements of $\mathbb Q[X,Y]$; repeated syntactic
representatives do not alter common zero sets. For a set $S$ of rational
points, the formal definition is exactly
\[
 \overline S=\{q:\ \forall f,\ (\forall r\in S,\ f(r)=0)\Rightarrow f(q)=0\}.
\]
This is the Zariski closure for rational-coefficient polynomials on rational
points. Quantification is over all polynomial expressions, not a selected
list of equations. Density means that this closure contains every point.
The coordinate formulas of $F$ are polynomial expressions; the code also
checks the iterate formula, the composition law for iterates, and a
polynomial representation of every iterate.

The final theorem \texttt{TLMC2685.counterexample} combines emptiness of the
multiplicative Jacobian zero set with failure of both alternatives. Its
second alternative explicitly asks for density of the orbit \emph{closure},
not merely of the original orbit. A separate theorem shows the closure is
not even contained in the preperiodic set. The Jacobian is computed from
recursively defined formal partial derivatives of the actual coordinate
polynomials.

The axiom audit reports only Lean's standard logical axioms
\texttt{propext}, \texttt{Classical.choice}, and \texttt{Quot.sound}.
There are no additional axioms, proof holes, \texttt{sorry}, or
\texttt{native\_decide} steps. The formal result deliberately makes no claim
that the orbit closure equals the line, or that the conjecture's unspecified
computability criterion has been independently defined in full generality.
Since the claimed dichotomy already fails, no analysis of that algorithmic
clause is needed for the counterexample.

\begin{thebibliography}{1}
\bibitem{tlmc}
The Last Math Competition, Conjecture 00000002685,
\href{https://github.com/The-Last-Math-Competition/The-Last-Math-Competition/blob/main/conjectures/00000002685.md}{repository source},
accessed 12 October 2026.
File blob SHA: \texttt{8d6a6cbadf48e960123f588074f14019237a7581}.
\end{thebibliography}
\end{document}
